\section{Discrete symmetries, the massless gauge algebra, and the states}
\label{sec:discrete}

\subsection{Charge conjugation, parity and time reversal}

At one constituent write $U=\langle e_2,e_3\rangle$ for the surplus plane, $\Omega=\gamma^0\gamma^1\gamma^2\gamma^3$, and
$\Gamma=J\Omega$. An operator $X$ on $S$ is \emph{linear} when $XJ=JX$ and \emph{antilinear} when $XJ=-JX$. It intertwines the free
Dirac operator $J\gamma^\mu\partial_\mu-m$ when it commutes with every $J\gamma^\mu$. Parity is required to be linear and time
reversal antilinear, and each must satisfy $XH(p)X^{-1}=H(-p)$ for every $p$ and $m$.

\begin{theorem}[Charge conjugation is the surplus plane]
\label{thm:C}
The antilinear operators that intertwine the Dirac operator form exactly the two-dimensional space $U$. Every unit vector $C\in U$
satisfies
\[
CH(p)C^{-1}=-H(-p),\qquad CJC^{-1}=-J,\qquad C\Gamma C^{-1}=-\Gamma,\qquad C\Lambda_+(p)C^{-1}=\Lambda_-(-p),
\]
commutes with every interior Lorentz generator, and preserves the mass form and every kinetic kernel $K^\mu$. Hence $C$ is unique up
to a gauge rotation $e^{\theta J}$, and the operator known as the Majorana real structure is charge conjugation.
\end{theorem}

\begin{proof}
If $XJ=-JX$, then $X(J\gamma^\mu)=(J\gamma^\mu)X$ is equivalent to $X\gamma^\mu=-\gamma^\mu X$. The conditions are linear, and
conjugation by a blade sends every blade to plus or minus itself, so the solution space is spanned by the blades that anticommute
with $J$ and with $e_1,e_4,e_5,e_6$. By \eqref{eq:bladecomm} a blade $e_I$ anticommutes with a vector $e_a$ exactly when
$|I|-[a\in I]$ is odd. For odd $|I|$ this forces $I\subset\{2,3\}$, hence $I=\{2\}$ or $\{3\}$, and both anticommute with $J$. For
even $|I|$ it forces $\{1,4,5,6\}\subset I$, hence $e_I=\Omega$ or $e_1\cdots e_6$, and both commute with $J$. The solutions are
therefore $e_2$ and $e_3$. The listed identities follow from the same rule, since $e_2,e_3$ anticommute with every $\gamma^\mu$, with
$J$ and with $\Gamma$, and commute with every interior bivector.
\end{proof}

\begin{corollary}[Equal mass, equal spin, opposite charge]
$C$ commutes with every generator of the interior Lorentz algebra, so it preserves the Casimir $-\sum_aS_a^2=\tfrac34\id$, and
$(CH(p)C^{-1})^2=E^2\id$. The Majorana condition $\Psi=C\Psi$ makes the current vanish: a particle that is its own antiparticle is
neutral.
\end{corollary}

\begin{theorem}[Parity and time reversal]
\label{thm:PT}
The admissible parities form the space $\langle\gamma^0,\gamma^0J\rangle$, and $P^2=+\id$ holds only at $P=\pm\gamma^0$, while $P=\pm\gamma^0J$
gives $P^2=-\id$. The admissible time reversals form $U\cdot\gamma^1\gamma^2\gamma^3$, so $T=C\gamma^1\gamma^2\gamma^3$ with
\[
T^2=-\id .
\]
The sign is not a convention: $B=\gamma^1\gamma^2\gamma^3$ has $B^2=+\id$ and anticommutes with every vector $C$ of the clock plane, so
$T^2=CBCB=-C^2B^2=-\|C\|^2\id$ for every choice of clock in the positive triple.
\end{theorem}

\begin{theorem}[$CPT$ lies on a circle]
\label{thm:CPT}
The product of the three operators depends on the parity chosen:
\[
P=\gamma^0:\quad CPT=-\Omega,\ (CPT)^2=-\id;\qquad\qquad P=\gamma^0J:\quad CPT=+\Gamma,\ (CPT)^2=+\id .
\]
What is convention-free is the orbit of $\Omega$ under $U(1)_J$, $\theta\mapsto\cos\theta\,\Omega+\sin\theta\,J\Omega$, whose four real points
$\pm\Omega,\pm\Gamma$ are the spacetime pseudoscalar and the chirality.
\end{theorem}

\begin{figure}[t]
\centering
\begin{tikzpicture}[scale=1.25,>=Stealth]
\draw[thick,gray] (0,0) circle (1.3);
\foreach \a/\l/\c in {0/{+\Omega}/scol,90/{+\Gamma=J\Omega}/tcol,180/{-\Omega}/scol,270/{-\Gamma}/tcol}{
  \fill[\c] (\a:1.3) circle (2pt);
  \node[\c] at (\a:1.85) {$\l$};
}
\draw[->,thick] (20:0.75) arc (20:70:0.75);
\node at (45:0.42) {\small $e^{\theta J}$};
\node[align=left,anchor=west] at (2.7,0.55) {\small $P=\gamma^0$: \quad $CPT=-\Omega$,\ $(CPT)^2=-\id$};
\node[align=left,anchor=west] at (2.7,-0.05) {\small $P=\gamma^0J$: \ $CPT=+\Gamma$,\ $(CPT)^2=+\id$};
\node[align=left,anchor=west] at (2.7,-0.65) {\small both lie on the $U(1)_J$ orbit of $\Omega$};
\end{tikzpicture}
\caption{The $CPT$ circle at one constituent (\cref{thm:CPT}). The product of the three discrete operators is the spacetime
pseudoscalar or the chirality according to the parity chosen; the two lie on one orbit of the gauge rotation $e^{\theta J}$.}
\label{fig:cpt}
\end{figure}

These are statements about explicitly constructed operators of the free theory; no $CPT$ theorem is proved, and the interacting
theory is not constructed.

\subsection{The massless gauge algebra at one constituent}

A gauge generator must leave the kinetic form invariant. Solving that condition over all of $\End_\R(S)$, with no ansatz, gives the
following.

\begin{proposition}[Gauge algebra from the kinetic forms]
\label{prop:gauge1}
The solutions of $X^{\!\top}K^\mu+K^\mu X=0$ for $\mu=0,\dots,3$ form
\[
\mathfrak g_{\rm massless}=\R\,\Omega\ \oplus\ \operatorname{span}\{J,\ e_2\Omega,\ e_3\Omega\}\ \cong\ \uu1_\Omega\oplus\su2 .
\]
The condition at $\mu=0$ is $X^{\!\top}=-X$, and the conditions at $\mu=a$ then say that $X$ commutes with the three boosts
$\gamma^0\gamma^a$ and hence with the whole interior Lorentz algebra. The trace form is $-8\,\id$ on this basis, so the algebra is
compact. Adding invariance of the mass form, $X^{\!\top}(\gamma^0J)+(\gamma^0J)X=0$, leaves exactly $\R J$. The algebra meets the span of the
pole operators in zero.
\end{proposition}

The $\su2$ factor contains the fermion-number direction $J$ and two generators that anticommute with $J$; it therefore mixes a field with
its conjugate and is the Pauli--G\"ursey algebra of a massless Dirac field \cite{Pauli57,Gursey58}. It is a symmetry of the free massless
theory, not a gauge candidate on this matter: its two Weyl components form a single doublet, and gauging $SU(2)$ on an odd number of
doublets is forbidden by the global anomaly \cite{Witten82}. The central factor is the axial one. Since $\Gamma=J\Omega$, $\Omega=-J\Gamma$,
which with $i=J$ and the chirality $\gamma_5=-\Gamma$ of \cref{thm:i} is the real form of $i\gamma_5$, as in \cref{prop:massangle}; $U(1)_\Omega$ is broken by a Dirac mass and is anomalous on this matter content \cite{Bilal,AlvarezGaumeWitten}.
The electromagnetic $\uu1$ is therefore singled out twice at one constituent: as the whole internal centraliser
(\cref{prop:6plus8plus1}) and as what survives of the massless gauge algebra once a mass is admitted.

\subsection{The state space}

\begin{theorem}[The real module is not a state space]
\label{thm:module}
Eight $16\times16$ integer matrices $\theta_i$, antisymmetric with $\theta_i^2=-\id$ and pairwise anticommuting, realise the canonical
anticommutation relations of the eight real field components exactly. Their commutant in $M_{16}(\R)$ is $\R\,\id$, because the $256$
Clifford monomials span the whole matrix algebra. Hence no internal complex structure commutes with all eight field operators, and the
quadratic observables built from them are real antisymmetric matrices with imaginary spectra. The quantum state space has to be formed
from outside, as $\mathcal F=\C\otimes_\R\R^{16}$, with charge $Q=-i\rho(J)$ and Hamiltonian $H=-i\rho(JH(p))$, where
$\rho(A)=-\tfrac14\sum_{ij}A_{ij}\theta_i\theta_j$ lifts an antisymmetric $8\times8$ matrix to the module.
\end{theorem}

It is the complex structure, not the reality, that is missing, which is the lesson of \cref{rem:ambient} arriving one level up.

\begin{theorem}[The free spectrum of one momentum]
\label{thm:spectra}
The lifted generators have characteristic polynomials
\[
\chi_{2\rho(J)}(x)=x^6(x^2+4)^4(x^2+16),\qquad \chi_{2\rho(JH)}(x)=x^6(x^2+4E^2)^4(x^2+16E^2),
\]
so $\spec Q=\{0,\pm1,\pm2\}$ and $\spec H=\{0,\pm E,\pm2E\}$, with multiplicities $6,4,4,1,1$. The joint eigenspaces of $(H,Q)$ are
\begin{center}
\begin{tabular}{@{}lll@{}}
\toprule
$H$ & dimension & charges\\
\midrule
$-2E$ & $1$ & $q=0$\\
$-E$ & $4$ & $q=\pm1$, twice each\\
$0$ & $6$ & $q=\pm2$ once each, $q=0$ four times\\
$+E$ & $4$ & $q=\pm1$, twice each\\
$+2E$ & $1$ & $q=0$\\
\bottomrule
\end{tabular}
\end{center}
so $H+2E$ has spectrum $0,E,2E,3E,4E$ with multiplicities $1,4,6,4,1$.
\end{theorem}

\begin{figure}[t]
\centering
\begin{tikzpicture}[xscale=1.2,yscale=0.95]
\foreach \y/\l in {-2/{-2E},-1/{-E},0/{0},1/{+E},2/{+2E}}{
  \draw[gray!60] (-2.6,\y) -- (2.6,\y);
  \node[left] at (-2.7,\y) {\small $\l$};
}
\node at (0,2.75) {\small charge $q$};
\foreach \q in {-2,-1,0,1,2}{\node at (\q,-2.7) {\small $\q$};}
% ground state
\fill[tcol] (0,-2) circle (2.4pt); \node[above right] at (0,-2) {\scriptsize 1};
% -E level
\fill[tcol] (-1,-1) circle (2.4pt); \fill[tcol] (1,-1) circle (2.4pt);
\node[above] at (-1,-1) {\scriptsize 2}; \node[above] at (1,-1) {\scriptsize 2};
% 0 level
\fill[tcol] (-2,0) circle (2.4pt); \fill[tcol] (2,0) circle (2.4pt); \fill[tcol] (0,0) circle (2.4pt);
\node[above] at (-2,0) {\scriptsize 1}; \node[above] at (2,0) {\scriptsize 1}; \node[above] at (0,0) {\scriptsize 4};
% +E
\fill[tcol] (-1,1) circle (2.4pt); \fill[tcol] (1,1) circle (2.4pt);
\node[above] at (-1,1) {\scriptsize 2}; \node[above] at (1,1) {\scriptsize 2};
% +2E
\fill[tcol] (0,2) circle (2.4pt); \node[above right] at (0,2) {\scriptsize 1};
\node[align=left,anchor=west] at (3.0,-2) {\small vacuum: unique, $q=0$};
\node[align=left,anchor=west] at (3.0,-1) {\small one particle or antiparticle};
\node[align=left,anchor=west] at (3.0,0) {\small two excitations};
\node[align=left,anchor=west] at (3.0,1) {\small three excitations};
\node[align=left,anchor=west] at (3.0,2) {\small filled};
\end{tikzpicture}
\caption{The Fock space of one momentum (\cref{thm:spectra}). Each dot is a joint eigenspace of energy and charge, labelled by its
dimension. The multiplicities $1,4,6,4,1$ are those of four fermionic modes: two particle and two antiparticle states of one Dirac
flavour, read off rather than inferred.}
\label{fig:fock}
\end{figure}

The multiplicities $1,4,6,4,1$ are the binomial coefficients of a four-mode Fock space, as a single Dirac flavour with two particle and
two antiparticle states must give. The ground state at $-2E$ is the zero-point energy $-\tfrac12\sum E$ of the four modes; summed over
momenta it is the familiar fermion term of the one-loop effective potential. A projector built from $\rho(JH)^2$ could not have
produced the table: squaring loses the sign of the energy and conflates the vacuum with the filled state.

\begin{proposition}[Conjugation on states]
\label{prop:Cstates}
$C\in SO(8)$ lifts to $\mathrm{Spin}(8)$ as a product of four integral reflections. The lift is unitary on $\mathcal F$, anticommutes with
$Q$ and carries $H(p)$ to $H(-p)$, so $\langle\psi|Q|\psi\rangle=0$ for every $C$-eigenstate. The stronger statement $Q|\psi\rangle=0$ is
false: the $+1$ eigenspace of $C$ has real dimension eight and $Q$ annihilates only a three-dimensional part of it.
\end{proposition}

\begin{remark}[Many-body chirality is a group lift]
$\rho$ is a homomorphism of Lie algebras only, so the product $\rho(J)\rho(\Omega)$, of rank eight, is not a $\mathrm{Spin}(8)$ element and
is not the chirality of the module. The correct object is $\Gamma_{\mathcal F}=\exp[\tfrac\pi2(\rho(J)+\rho(\Omega))]$, an involution of
$\mathcal F$ whose eigenvalue on a joint eigenvector is $i^{q_J+q_\Omega}$; the sectors are labelled by the parity of
$\tfrac12(q_J+q_\Omega)$.
\end{remark}
